% ============================================================
% CHAPTER 9: CONCLUSION
% Citations reuse verified entries in bibliography.bib only.
% ============================================================

This thesis asked whether the dynamical vocabulary of complex systems can be
applied empirically to LLM-based multi-agent debate, and what such a treatment
reveals that outcome metrics cannot. To answer it we built a two phase debate
protocol whose logs expose private beliefs and confidences alongside public
messages, defined \suitename{} as four motif detectors each grounded in the
complex systems or opinion dynamics literature and each tested against its own
null model, and ran the full grid on two deliberately contrasting benchmarks
across $30{,}000$ debates.

The first research question asked whether the motifs can be measured at all.
\suitename{} detects them, at very different prevalences, and the differences are
themselves informative. Self-reinforcement is universal, with a prevalence
between $0.733$ and $0.881$ in every combination of dataset, topology, and memory
window. Multistability is common and belongs to the question rather than to the
configuration, since $76\,\%$ of the $600$ task and configuration cells admit more
than one stable outcome and per-task values are consistent across configurations.
Dominance emerges behaviourally even where the wiring is symmetric, concentrating
above the topology-keyed null in all but one of the hundred task and benchmark
combinations. Limit cycles are rare and real, since almost every sequence reaches
a fixed point and the small residue yields $47$ agent level and $32$ system level
flagged cases, most of them clean period two orbits on inspection. That they sit
almost entirely in the star topology, and that the fully connected topology
produces none, is a clean result about cooperative protocol design rather than a
measurement failure.

The second research question asked what governs the motifs and what they predict.
The two structural axes shape the dynamics and not the correctness of the debate.
The star topology is slower on both benchmarks, concentrates influence by
construction, and is the only source of genuine cycles, while the memory window
affects convergence speed only and multistability is invariant to both. What
governs the outcome instead is the composition of the opening votes.
Self-reinforcement is the convergence mechanism everywhere, yet it strengthens
whichever option leads at the outset rather than the correct one, so faster
convergence is not more accurate convergence. Multistability decides how many
outcomes are available to be selected and helps only where the single basin the
group would otherwise fall into is the wrong one. Dominance helps when the debate
elects its own hub, because an elected hub started correct far more often than a
random agent, and hurts when the hub is prescribed by the wiring, because a fixed
centre spreads the shared bias instead of correcting it. Limit cycles mark pure
efficiency failures, since every one of them ran to the cap without producing a
decision.

The third research question asked whether the opening votes are causally
responsible for the outcome. They are. Observationally, whether the ground truth
appears at all in the opening distribution is the strongest predictor of accuracy
we measured. Causally, seeding one correct agent into an otherwise wrong opening
raises accuracy from $1.8\,\%$ to $26.2\,\%$ and the effect holds in every task
tested, while the system almost never recovers an answer no agent initially held.
The same experiment shows the limits of the effect, since a correct supermajority
that a plurality vote would preserve with certainty is talked out of the right
answer in roughly one repetition in seven, and a single agent that never updates
collapses accuracy from $63.7\,\%$ to $10.0\,\%$ even when outnumbered two to one
by correct agents. Stated with its scope, the overarching finding is that within
a homogeneous, cooperative, four agent debate on one model family across two
benchmarks, deliberation works predominantly as a selective amplifier of its
initial vote distribution rather than as a mechanism that reasons a new answer
into existence.

The three contributions therefore close on each other. The instrument made the
characterisation possible, because motif prevalence is only meaningful against a
null and attractor statistics are only visible under repeated runs of the same
task. The characterisation made the causal test necessary, because three of the
four motifs pointed back to the same opening distribution and only an
intervention could separate that from task difficulty. The causal test in turn
explains what the instrument was measuring all along, since each detector reads
one face of a single amplification process. That is also why \suitename{}
survives the deflationary result. If the debate mostly magnifies what it started
with, then the quantities worth logging are no longer accuracy and confidence but
the shape of the opening distribution and the concentration of influence that
follows from it, and those are exactly the quantities these four detectors
compute.

Multi-agent debate is often presented as a way to make language models reason
better together \citep{du2024improving, liang2024encouraging}. Our results
suggest a more precise and more cautious picture, in which the group mostly
magnifies what its members already believe, so the interesting question is not
whether the debate converged but on what it converged and from where.
\suitename{} turns that question into four measurements computable from ordinary
interaction logs, and its broader value is as groundwork for a diagnostic
instrument that would let a practitioner check which dynamical regime a
deliberation is in before trusting what it produces.